\documentclass[tikz,border=8pt]{standalone}

\usepackage{amsmath}
\usepackage{bm}
\usepackage{tikz-3dplot}
\usetikzlibrary{arrows.meta,calc}

\begin{document}

\tdplotsetmaincoords{65}{115}

\begin{tikzpicture}[
    line cap=round,
    line join=round,
    >=Stealth,
    font=\small
]

\def\R{2.35}

\tikzset{
    spheregrid/.style={
        draw=gray!28,
        line width=.55pt
    },
    sphereedge/.style={
        draw=gray!58,
        line width=1.0pt
    },
    hiddenedge/.style={
        draw=gray!32,
        dashed,
        dash pattern=on 2pt off 2.3pt,
        line width=.65pt
    },
    ray/.style={
        -{Stealth[length=6pt,width=6pt]},
        draw=black!70,
        line width=.9pt
    },
    tangent/.style={
        -{Stealth[length=5pt,width=5pt]},
        draw=gray!58,
        line width=.75pt
    },
    patch/.style={
        draw=black!62,
        line width=.75pt,
        fill=black!45,
        fill opacity=.18
    },
    surfpoint/.style={
        circle,
        fill=black!70,
        inner sep=1.05pt
    },
    centerpoint/.style={
        circle,
        fill=black!65,
        inner sep=.85pt
    },
    labelbox/.style={
        fill=white,
        fill opacity=.82,
        text opacity=1,
        rounded corners=1.6pt,
        inner xsep=2.6pt,
        inner ysep=1.4pt,
        align=center
    }
}

% =========================================================
% Draw an upper hemisphere z >= 0 using true 3D parametric curves
% =========================================================
\newcommand{\drawhemi}{%
    % equator: back half dashed, front half solid
    \draw[hiddenedge]
        plot[domain=25:205, samples=80, variable=\t]
        ({\R*cos(\t)}, {\R*sin(\t)}, 0);

    \draw[sphereedge]
        plot[domain=205:385, samples=80, variable=\t]
        ({\R*cos(\t)}, {\R*sin(\t)}, 0);

    % latitude circles
    \foreach \th in {24,46,68}{
        \pgfmathsetmacro{\rr}{\R*sin(\th)}
        \pgfmathsetmacro{\zz}{\R*cos(\th)}
        \draw[spheregrid]
            plot[domain=25:385, samples=100, variable=\t]
            ({\rr*cos(\t)}, {\rr*sin(\t)}, \zz);
    }

    % meridians
    \foreach \ph in {25,70,115,160,205,250,295,340}{
        \draw[spheregrid]
            plot[domain=0:90, samples=70, variable=\t]
            ({\R*sin(\t)*cos(\ph)},
             {\R*sin(\t)*sin(\ph)},
             {\R*cos(\t)});
    }

    % two stronger silhouette-like meridians
    \draw[sphereedge]
        plot[domain=0:90, samples=70, variable=\t]
        ({\R*sin(\t)*cos(25)},
         {\R*sin(\t)*sin(25)},
         {\R*cos(\t)});

    \draw[sphereedge]
        plot[domain=0:90, samples=70, variable=\t]
        ({\R*sin(\t)*cos(205)},
         {\R*sin(\t)*sin(205)},
         {\R*cos(\t)});
}

\begin{scope}[tdplot_main_coords]
    \coordinate (O) at (0,0,0);

    % Hemisphere: no gradient, no ball shading.
    \drawhemi

    % ---------------------------------------------------------
    % Local surface point and local solid-angle element
    % The patch lies on x^2 + y^2 + z^2 = R^2, z >= 0.
    % ---------------------------------------------------------
    \pgfmathsetmacro{\px}{-0.50}
    \pgfmathsetmacro{\py}{0.80}
    \pgfmathsetmacro{\hx}{0.20}
    \pgfmathsetmacro{\hy}{0.16}

    \pgfmathsetmacro{\xA}{\px-\hx}
    \pgfmathsetmacro{\xB}{\px+\hx}
    \pgfmathsetmacro{\yA}{\py-\hy}
    \pgfmathsetmacro{\yB}{\py+\hy}

    \pgfmathsetmacro{\pz}{sqrt(\R*\R-\px*\px-\py*\py)}
    \coordinate (P) at (\px,\py,\pz);

    % Curved local patch on the hemisphere.
    \path[patch]
        plot[domain=\xA:\xB, samples=14, variable=\t]
            ({\t},{\yA},{sqrt(\R*\R-\t*\t-\yA*\yA)})
        -- plot[domain=\yA:\yB, samples=14, variable=\t]
            ({\xB},{\t},{sqrt(\R*\R-\xB*\xB-\t*\t)})
        -- plot[domain=\xB:\xA, samples=14, variable=\t]
            ({\t},{\yB},{sqrt(\R*\R-\t*\t-\yB*\yB)})
        -- plot[domain=\yB:\yA, samples=14, variable=\t]
            ({\xA},{\t},{sqrt(\R*\R-\xA*\xA-\t*\t)})
        -- cycle;

    % Direction vector omega_o.
    \draw[ray] (O) -- (P);

    % Tangent directions at P:
    % z = sqrt(R^2 - x^2 - y^2)
    % tangent basis:
    % e_x = (1,0,-x/z), e_y = (0,1,-y/z)
    \pgfmathsetmacro{\sone}{0.48}
    \pgfmathsetmacro{\stwo}{0.46}

    \coordinate (Tone) at
        ({\px+\sone},
         {\py},
         {\pz-\sone*\px/\pz});

    \coordinate (Ttwo) at
        ({\px},
         {\py+\stwo},
         {\pz-\stwo*\py/\pz});

    \draw[tangent] (P) -- (Tone);
    \draw[tangent] (P) -- (Ttwo);

    \node[centerpoint] at (O) {};
    \node[surfpoint] at (P) {};
\end{scope}

% ------------------------------------------------------------
% Labels: all black, with translucent white background.
% ------------------------------------------------------------
\node[labelbox, text=black] at ($(O)!0.55!(P)+(0.10,-0.02)$)
    {$\boldsymbol{\omega}_o$};

\node[labelbox, text=black, anchor=east] at ($(Tone)+(-0.08,0.09)$)
    {$\boldsymbol{\omega}_1$};

\node[labelbox, text=black, anchor=west] at ($(Ttwo)+(0.18,0.03)$)
    {$\boldsymbol{\omega}_2$};

\node[labelbox, text=black, anchor=west] at ($(P)+(0.50,0.22)$)
    {$\mathrm{d}\boldsymbol{\omega}_o$};

\node[text=black!75] at (0,-1.02) {$S^2$};

\end{tikzpicture}

\end{document}